\scanpage{21}
% L1–2
Soit $\unsure{\mathbb{\Delta}|\Delta}$ la catégorie des ens finis $\mathrm{tot}^{\mathrm t}$ ordonnés \inserted{non vides}, avec comme morphismes les applications \uline{strictement} croissantes.
% L3
Soit
\[
\unsure{\mathbb D|D}_{*}\in\mathrm{Ob}\widehat{\mathbb{\Delta}}
\qquad \mathbb D_*:\mathbb{\Delta}^{\circ}\longrightarrow(\mathrm{Ens})
\]
% L4–7
donc $\mathbb D_*$ définit une catégorie fibrée à fibres discrètes sur $\mathbb{\Delta}$, à quelle condition celle-ci est-elle ordonnée \inserted{(ou préordonnée, celà revient au $=$)} ? Posons $\Delta_n=[0,n]\in\mathrm{Ob}\,\mathbb{\Delta}$, $\mathbb D_n=\mathbb D_*(\Delta_n)$, la condition est celle-ci
% L8–10
\textcircled{1} $\forall n,m\in\mathbb N$, avec $n<m$, et $d_m\in\mathbb D_m$, l'application
\[
\mathrm{Hom}_{\mathbb{\Delta}}(\Delta_n,\Delta_m)\longrightarrow\mathbb D_n
\qquad u\longmapsto u^*(d_m)
\]
est injective.

% L11–13
Ceci posé, on a donc un ens ordonné $\mathbb D_*=\coprod_{m\in\mathbb N}D_m$ et un foncteur fibrant
\[
\underline{D}_*\longrightarrow\mathbb{\Delta}
\]
% L14–18
Considérons la réalisation topologique \inserted{$|D_*|$} de $D_*$, donc $\forall d\in D_*=\mathrm{Ob}\,\underline D$, on a une partie \inserted{fermée} $|d|\subset|D_*|$, l'application $d\longmapsto|d|$ de
\[
D_*\longrightarrow\mathscr P_f(|D_*|)
\]
étant injective et préservant la relation d'ordre.
% L19–21
Considérons \inserted{pour $d_m\in\mathbb D_m$ \deleted{\unsure{$m$}}} l'ens $D_d=\{\delta\in D\mid\delta\leq d\}$, on trouve qu'il est $\mathrm{can}^{\mathrm t}$ isom (comme ens. ordonné) à $\mathscr P^*(\Delta_m)$ ($\mathscr P^*$ désigne l'ens des parties $\ne\varnothing$).
% L22–24
Donc $|d|$ est canoniquement isomorphe au simplexe type $|\Delta_m|$. Ainsi on trouve
\[
|D_*|=\varinjlim_{d\in D_*}|d|=\varinjlim_{d\in D_*}|\Delta_m|
=\varinjlim_{\Delta\in\mathbb{\Delta}/D}|\Delta|
\]
% archive sheet number: 20
